\section{Results and Discussion}
\label{sec:results}
\subsection{In-Distribution: Three Seeds}
On identical tests, the base recipe gives typed-decisions 55.9/53.3/55.6
(mean 54.9), BANKING77 76.8/83.4/83.8 (mean 81.3), and CLINC150
68.2/76.6/73.8 (mean 72.9). Seed spread on BANKING77 (76.8 vs $\sim$83.5) is
exactly why single-seed reporting is insufficient. Untrained-head baselines sit
at chance (31\%/2\%/17\%).

\begin{table}[htbp]
\caption{In-distribution accuracy on identical tests (500/500/400 items).}
\label{tab:seeds}
\begin{center}
\setlength{\tabcolsep}{4pt}
\begin{small}
\begin{tabular}{lccc}
\toprule
Seed & typed & BANKING77 & CLINC150 \\
\midrule
3407 & 55.9 & 76.8 & 68.2 \\
41 & 53.3 & 83.4 & 76.6 \\
42 & 55.6 & 83.8 & 73.8 \\
Mean & 54.9 & 81.3 & 72.9 \\
\bottomrule
\end{tabular}
\end{small}
\end{center}
\end{table}

\subsection{Transfer: Leave-One-Taxonomy-Out}
With BANKING77 fully held out, accuracy reaches \textbf{44\% [38--47\%]}
against 2\% chance (paraphrase median 42.8\%, range 42.2--43.6); with CLINC150
held out, \textbf{55.6\%} (paraphrase 55.6--57.0) against 17\% chance. Transfer
is substantial but partial versus 81.3\%/72.9\% in-distribution.

Transfer is wording-robust, not prompt luck: a 3-variant grid (instruction
rewording plus option-order shuffling) moves LOTO BANKING77 only across
42.2--43.6\% (median 42.8\%) and LOTO CLINC150 across 55.6--57.0\%, while
in-distribution A holds 86.0--86.4\% on BANKING77. Ranges this tight mean the
effect lives in schema semantics, not phrasing.

\begin{table}[htbp]
\caption{Paraphrase grid: median [range] over 3 wording/option-order variants.}
\label{tab:para}
\begin{center}
\setlength{\tabcolsep}{3pt}
\begin{footnotesize}
\begin{tabular}{lcc}
\toprule
Model & BANK77 & CLINC150 \\
\midrule
LOTO banking & 42.8 [42.2--43.6] & --- \\
LOTO clinc & --- & 56.0 [55.6--57.0] \\
A & 86.0 [86.0--86.4] & 84.6 [84.4--84.8] \\
E & 79.4 [79.0--80.2] & 70.0 [69.8--70.4] \\
\bottomrule
\end{tabular}
\end{footnotesize}
\end{center}
\end{table}

\begin{table}[htbp]
\caption{LOTO transfer vs in-distribution means (accuracy \%).}
\label{tab:loto}
\begin{center}
\setlength{\tabcolsep}{4pt}
\begin{small}
\begin{tabular}{lcc}
\toprule
Setting & BANKING77 & CLINC150 \\
\midrule
Held out (transfer) & 44 [38--47] & 55.6 \\
In-distribution mean & 81.3 & 72.9 \\
Chance & 2 & 17 \\
\bottomrule
\end{tabular}
\end{small}
\end{center}
\end{table}

\subsection{Ablation: The Truncation Confound}
Row-matched canonical (A) vs full-augmentation (E) training gives:

\begin{table}[htbp]
\caption{Ablation A vs E (28\,000 rows each; decisions 31\,754 vs 39\,560).}
\label{tab:ablation}
\begin{center}
\setlength{\tabcolsep}{4pt}
\begin{small}
\begin{tabular}{lccc}
\toprule
Treatment & typed & BANKING77 & CLINC150 \\
\midrule
A canonical & 54\% [51--56] & \textbf{86\%} [81--88] & \textbf{84\%} [81--88] \\
E augmented & 59\% [56--61] & 79\% [75--82] & 70\% [65--73] \\
\bottomrule
\end{tabular}
\end{small}
\end{center}
\end{table}

An earlier 36\%\,$\to$\,83\% jump attributed to augmentation was a context
truncation confound: training at 2\,048 tokens cut wide schemas, and
augmentation's option subsampling merely made them fit. With fitting context,
canonical training wins in-distribution; augmentation buys transfer plus memory
savings (full 77/150-option schemas OOM at batch 8 on T4) at a small in-domain
tax. Residual confounds disclosed: +25\% decisions for E (biasing toward E) and
batch 4 vs 8 at equal effective size 32.

Overfitting dynamics are strikingly consistent: eval-loss falls through epoch 2
then spikes at epoch 3 in every fully-logged run, so best-checkpoint selection
always restores epoch 2 (Table~\ref{tab:overfit}).

\begin{table}[htbp]
\caption{Eval-loss at epochs 2 and 3 (best checkpoint underlined).}
\label{tab:overfit}
\begin{center}
\setlength{\tabcolsep}{4pt}
\begin{small}
\begin{tabular}{lcc}
\toprule
Run & Epoch 2 & Epoch 3 \\
\midrule
seed41 & \underline{0.72} & 0.90 \\
seed42 & \underline{0.62} & 1.01 \\
LOTO-banking & \underline{0.59} & 0.72 \\
LOTO-clinc & \underline{0.64} & 0.87 \\
Ablation A & \underline{0.76} & 0.95 \\
Ablation E & \underline{0.57} & 0.65 \\
\bottomrule
\end{tabular}
\end{small}
\end{center}
\end{table}

\subsection{Calibration Under Shift}
On BANKING77, ECE rises from \textbf{0.026} in-distribution to \textbf{0.366}
held-out (Brier 0.281\,$\to$\,0.876); autonomous coverage at 5\% risk collapses
from 63\% to 0.6\%. Temperature scaling fit in-distribution does not transfer.
Yet error-detection AUROC falls only from 0.87 to 0.73: confidence still
separates right from wrong answers, confirming that accuracy, calibration, and
sensitivity are distinct quantities~\cite{systemswitch2026}. Measured T4 latency
is 56--249\,ms/item depending on schema width.

\begin{table}[htbp]
\caption{Calibration per bench: LOTO (shift) vs E (in-distribution).}
\label{tab:rq3}
\begin{center}
\setlength{\tabcolsep}{3pt}
\begin{footnotesize}
\begin{tabular}{llcccc}
\toprule
Model & Bench & Acc & ECE & AUROC & Cov@5\% \\
\midrule
LOTO & typed & 59 & 0.053 & 0.64 & 3 \\
LOTO & banking* & 43 & 0.366 & 0.73 & 0.6 \\
LOTO & clinc & 65 & 0.176 & 0.82 & 29 \\
E & typed & 59 & 0.121 & 0.62 & 0.1 \\
E & banking & 80 & 0.026 & 0.87 & 63 \\
E & clinc & 70 & 0.090 & 0.85 & 40 \\
\bottomrule
\end{tabular}
\end{footnotesize}
\end{center}
\end{table}

\subsection{External Context}
Table~\ref{tab:external} places our accuracy beside vendor-reported numbers.
Metric and protocol boundaries differ per row (Decision Index reports macro-F1
over full tests; Unsloth reports accuracy under its own mixture; ours is
accuracy on identical 500/500/400 tests), so no cross-row ranking is claimed.
Two observations survive the boundary: (i)~a deployed 421M system (Laya) sits
near chance on intent taxonomies while our 350M reaches 81--86\%, i.e.\ recipe
dominates size in this class; (ii)~9--27B systems lead absolutely (90--97\%),
as expected, with Cloudflare's own rows flagged self-reported by the board.

\begin{table}[htbp]
\caption{Context table. Ours: accuracy on identical tests. Starred rows:
vendor-reported under own protocols and metrics (acc vs macro-F1).}
\label{tab:external}
\begin{center}
\setlength{\tabcolsep}{4pt}
\begin{small}
\begin{tabular}{llccc}
\toprule
Model & Params & BANK77 & CLINC150 & typed \\
\midrule
Ours mean (acc) & 350M & 81.3 & 72.9 & 54.9 \\
Ours canonical A (acc) & 350M & 86 & 84 & 54 \\
Qwen3.5-0.8B* (acc) & 0.8B & 74 & 76 & 73 \\
Clef* (F1) & 27B & 94.2 & 97.4 & --- \\
Clef-flash* (F1) & 9B & 90.9 & 66.8 & --- \\
Jev* (F1) & closed & 79.7 & 89.3 & --- \\
Kev 9B* (F1) & 9B & 84.8 & 79.0 & --- \\
Laya* (F1) & 421M & 14.3 & 3.2 & --- \\
\bottomrule
\end{tabular}
\end{small}
\end{center}
\end{table}

\subsection{Capacity Frontier}
Typed-decisions plateaus at 54--59\% across all seven runs while Qwen3.5-0.8B
reports 73\%: unlike single wide taxonomies (handled well even at 150
classes), its records demand up to five interdependent judgments per state
(risk, review need, routing) that stretch 16 layers at width 1\,024. The gap
is architectural capacity, not recipe, and it bounds deployment scope to
single-question routing rather than joint multi-field reasoning.

\subsection{Discussion}
Fixed 350M capacity forces mixture trade-offs: doubling the typed-decisions
share shifted BANKING77 +2 but CLINC150 $-$7 without lifting either, i.e.\
weights relocate capacity rather than create it. For practitioners: fit schemas
in context first (truncation dwarfs recipe effects), buy augmentation for transfer and memory rather than in-domain
points, and never trust in-distribution temperatures under shift. Deploy with
an escalation policy
\begin{equation}
\hat{y}(x) = \arg\max_k p_k(x) \text{ if } \max_k p_k(x)\ge\tau
\text{ else ESCALATE},
\end{equation}
with $\tau$ set from AUROC-validated coverage analysis (our models sustain 63\%
autonomy in-distribution at 5\% risk but near zero under shift, so $\tau$ must
be re-fit per deployment distribution). Limits: A/E are single-seed
with a batch-size confound (4$\times$8 vs 8$\times$4 at equal effective size 32,
forced by full-schema OOM); rename/paraphrase/derived types are not yet
separated, so the transfer credit cannot be attributed to one mechanism; no
same-protocol Qwen rerun exists (vendor numbers are labeled context only);
OOS/OOD slices and macro-F1 leaderboard slices are future work; and the
\texttt{lfm1.0} license text is still unread, which must precede any release
claim. None of these touch the three headline findings, each of which rests on
fixed test sets with disclosed residuals.
